% D9, after D1 (collective acceptance) to D8 (defeasibility).
\subsection{Seals relocate authority without constraining it}\label{sec:obj-relocate}

\paragraph{The objection.} A signature shows who issued a record, and a validator shows that the
record met the rules of its kind. Neither detects a decision made in bad faith by someone acting
within valid standing. Trust in this design is indexed by its observer, and nothing enforces a
verdict on everyone, so the sort moves authority to whoever holds a seal and leaves it unchecked
there. A design that spreads records across many hands without governance, dispute resolution or
consequences for abuse could then turn out worse than the central systems it means to replace.

\paragraph{Reply.} I accept the premise. The thesis concerns the necessity of authority, fact by
fact, and leaves open which authority is legitimate. \citet{reijers2016} put the second question to
blockchain governance; this paper puts the first, and a seal does nothing on its own to make its
holder just.

The objection still overreaches in its conclusion, for reasons \cref{sec:anchoring} develops. A
charge of abuse comes apart along the lines of the sort: its procedural parts are monotone and
anyone holding the records can check them, its parts that ask for an absence need a seal, and its
parts about intent need judgment, which the sort locates without performing (\cref{sec:charge}).
The objection has the shape of Hart's case against a regime of primary rules alone, and the remedy
takes the shape of his: rules of recognition, change and adjudication, written as records and
co-signed by the parties they bind (\cref{sec:secondary}). With those rules in place, the comparison
with a central system favors the design of records, since a platform seals in private while every
seal here is a signed record that names its signers (\cref{sec:central}).

The concern also has independent support from engineering. Builders of local-first group
management met the same problem as duelling administrators who revoke each other concurrently, and
concluded that an arbiter outside the merge has to order concurrent revocations, because a Byzantine
administrator can otherwise exploit the concurrency \citep{dougal2026}. The sort predicts exactly
this: revocation asks for an absence, so it needs a sealed order (\cref{sec:suspension}).

What the sort cannot supply is the content of the secondary rules: who sits on a panel, which
quorum may suspend an office, how long a suspension takes to reach every reader. Those are choices a
group makes when it anchors its frame. The paper claims to show where such choices have to be made,
and that a design of records can make them in view of everyone they bind.
